% ============================================================
% Sec. 2 draft: why the two objectives select different units
% Numbers: paper-analysis/2026-09-24_finding1_numbers.md (scripts/finding1_numbers.py).
% Shares: per-layer additive shares, averaged over the layers of each band;
% text = prompt + generated CoC. Rank agreement: raw (uncorrected) Spearman.
% Bands: layers 0-21 ("early and middle") and 22-34 ("late").
% Main text and figures use Q heads; MLP channels go to the supplement.
% Terms: "CoC score" = first-order Taylor score under the CE loss of the CoC;
% "trajectory score" = first-order Taylor score under the FM loss.
% TODO: use the same two terms in Sec. 1 and the method section.
% ============================================================

% TODO: renumber findings (Sec. 1 is Finding 1; this is Finding 2).

\subsection{Why the Two Objectives Select Different Units}
\label{sec:vlm_analysis}

\noindent\textbf{Finding 2.}\ \emph{The CoC and trajectory scores rank units
alike in the early and middle layers but differently in the late layers, because
there the two scores come from different token positions.}

\paragraph{The rankings part in the late layers.}
For each unit we compute a first-order Taylor score under each of the two
losses: the cross-entropy of the generated CoC (the \emph{CoC score}) and the
flow-matching loss of the trajectory (the \emph{trajectory score}).
Within each VLM layer, we rank Q heads by each score and compare the two
rankings with Spearman's $\rho$ (Figure~\ref{fig:vlm_analysis}a).
In layers 0--21 the rankings largely agree ($\rho=0.78$). In layers 22--34 they
part ($\rho=0.33$).
The two scores thus favor similar heads in the early and middle layers and
different heads in the late layers.
MLP channels show the same drop ($0.75\to0.34$), as do all results below
(supplement).

\paragraph{Where each score comes from.}
A unit produces an output at every token position, and its Taylor score is a sum
of the contributions of these outputs. Grouping the contributions by the kind of
token splits each score into shares from vision, ego-history, prompt-text and
generated-CoC positions (Figure~\ref{fig:vlm_analysis}b,c).
% NOTE: exact definition (signed per-type shares that sum to the clip-level
% |sum_p a_p g_p|) goes in the appendix if needed.

In layers 0--21, most of the trajectory score comes from vision positions
(58\%). The CoC score is split between vision (47\%) and text positions (48\%),
a text share nearly twice that of the trajectory score (26\%). This difference
in text share is the main difference between the two scores in these layers.

In layers 22--34 the two scores part. The CoC score comes almost entirely from
text positions (96\%). The trajectory score stays spread over vision (34\%),
text (53\%) and ego-history positions (13\%), and most of its text share comes
from the generated CoC (33 of 53 points).
In the late layers, the CoC score therefore favors heads that process language
positions, whereas the trajectory score favors heads that also read the scene,
the ego history and the generated reasoning.

\paragraph{Position explains most of the disagreement.}
When both scores are restricted to the same kind of position, the late-layer
rankings agree again: $\rho$ rises from $0.33$ to $0.72$ at vision, $0.75$ at
prompt-text and $0.86$ at ego-history positions (Figure~\ref{fig:vlm_analysis}d).
The exception is the generated-CoC positions ($0.36$): even at the same tokens,
the two scores favor different heads.
% NOTE (MLP, supplement): 0.34 -> 0.47-0.52 at vision / prompt-text /
% ego-history, 0.24 at generated-CoC positions. MLP channel scores are noisier
% (split-half reliability 0.46-0.49 in layers 22-34 vs. 0.81-0.89 for Q heads);
% corrected for this, MLP rises from 0.54 to 0.71-0.83 (CoC positions 0.38).

% ------------------------------------------------------------
% Figures
% ------------------------------------------------------------
% fig:vlm_analysis (main, one figure, four panels, Q heads; one row of four,
% or 2x2 with (a, d) on top and (b, c) below)
%   (a) per-layer Spearman rho between CoC and trajectory rankings of Q heads,
%       layers 0-34 (fig3_vlm_rank with the MLP line removed). The Q-head curve
%       oscillates in the late layers (32 heads per layer); overlay the 0-21 and
%       22-34 band means as horizontal lines, or add a per-layer bootstrap band.
%   (b) per-layer shares of the trajectory score at vision / text / ego-history
%       positions, Q heads.  TODO: generate (existing fig3_vlm_traj_share is MLP?).
%   (c) per-layer shares of the CoC score, Q heads.  TODO: generate.
%   (d) layers 22-34, rank agreement with both scores restricted to one kind of
%       position: bars for all positions / vision / prompt-text / ego-history /
%       generated-CoC (0.33 / 0.72 / 0.75 / 0.86 / 0.36). TODO: generate.
%   Caption draft: "CoC and trajectory scores of Q heads in the VLM.
%   (a) Agreement between the two rankings within each layer. (b, c) Share of
%   each layer's trajectory score (b) and CoC score (c) that comes from each kind
%   of token position; late layers (22-34) are shaded. (d) Agreement in layers
%   22-34 when both scores are restricted to one kind of token position."
% Supplement:
%   - MLP versions of (a)-(d).
